\documentclass[11pt]{article}
\usepackage{amsmath,amssymb,amsthm}
\newtheorem{lemma}{Lemma}
\newtheorem{theorem}{Theorem}
\title{Geometric Cauchy localization and the full-cube slot budget}
\author{}\date{}
\begin{document}\maketitle

\paragraph{Status and scope.}
This is a complete written obstruction to one exact bridge, accompanied
by exact rational checks. Geometric Cauchy kernels do localize well on
exterior queries. However, a fixed head slot which can receive every
geometric pole token cannot preserve that kernel, even up to separable
amplitudes, inside a full-cube-admissible reverse denominator. Three
tokens suffice for the obstruction. This is not a lower bound for sums
of admissible heads or for sign-table complexity. It does not exclude
target-adaptive routing, arbitrary affine numerators, or joint
cancellations among several heads. It has not been formalized in Lean.

\begin{lemma}[A three-index budget constraint]
Suppose
\[
 B_{ij}=b-p_i-q_j,\qquad p_i,q_j\geq0,\qquad
 \sum_i p_i+\sum_jq_j<b.                                      \tag{1}
\]
For any three distinct matched row and column indices, there are two
of them, denoted $i,j$, for which
\[
 \frac{B_{ii}B_{jj}}{B_{ij}B_{ji}}>\frac14.                    \tag{2}
\]
The conclusion also holds for a selected three by three submatrix of
a larger matrix satisfying (1).
\end{lemma}
\begin{proof}
All denominators are positive and at most $b$. Among three nonnegative
numbers $p_i+q_i$ whose sum is strictly smaller than $b$, at least two
are strictly smaller than $b/2$. Their diagonal denominators exceed
$b/2$, while $B_{ij}B_{ji}\leq b^2$. This proves (2).
Restricting the sum in (1) to selected rows and columns preserves the
strict inequality.
\end{proof}

For a matrix with nonzero entries, define its inverse diagonal
cross-ratio on indices $i,j$ by
\[
 \mathcal R_{ij}(K)=\frac{K_{ij}K_{ji}}{K_{ii}K_{jj}}.          \tag{3}
\]
This quantity is unchanged on replacing $K_{ij}$ by $u_i v_j K_{ij}$
for arbitrary nonzero separable row and column factors. In particular,
for $K_{ij}=1/B_{ij}$ it equals the ratio in (2).

\begin{theorem}[No exact three-scale bridge]
Let $0<a_1<a_2<a_3$ with $a_{i+1}/a_i\geq16$, and put
$C_{ij}=1/(a_i+a_j)$. There do not exist a denominator matrix satisfying
(1) and nonzero row and column factors $u_i,v_j$ such that
\[
                    \frac1{B_{ij}}=u_i v_j C_{ij}
                       \quad(1\leq i,j\leq3).                \tag{4}
\]
More generally, let $L\geq1$. If every successive scale ratio is at
least $16L^4$, then there is no such representation with multiplicative
errors $E_{ij}\in[L^{-1},L]$ on the right side of (4).
\end{theorem}
\begin{proof}
For $a_j/a_i=r\geq16$,
\[
 \mathcal R_{ij}(C)=\frac{4a_i a_j}{(a_i+a_j)^2}
                   =\frac{4r}{(1+r)^2}<\frac4r\leq\frac14.
                                                               \tag{5}
\]
Every distinct pair has this property, contradicting (2) and the
invariance of (3). With multiplicative errors, the ratio in (3)
changes by the factor $E_{ij}E_{ji}/(E_{ii}E_{jj})\leq L^4$.
If $r\geq16L^4$, it remains strictly smaller than $1/4$.
\end{proof}

\paragraph{Why fixed slotwise projective changes are covered.}
For a nonconstant fractional linear map
$\phi(x)=(\alpha x+\beta)/(\gamma x+\delta)$,
\[
 \phi(v)-\phi(u)
 =\frac{(\alpha\delta-\beta\gamma)(v-u)}
        {(\gamma v+\delta)(\gamma u+\delta)}.                 \tag{6}
\]
Away from the pole of $\phi$, the reciprocal difference therefore
changes only by separable factors. Taking $u=-a_i$ and $v=a_j$
shows that a fixed projective change for one head slot, together with
arbitrary nonzero separable amplitudes, is included in the theorem.
If every one of three tokens can be routed to this slot across input
rows and their three matching queries are present, the corresponding
submatrix already contradicts full-cube admissibility. Additional
rows, columns, and zero-residue padding do not improve its denominator
budget. Zero factors cannot reproduce the nonzero entries of $C$.
If a proposed representation instead discards these entries or uses
their cancellation against other heads, this theorem does not apply.

\paragraph{Localization before enforcing the cube budget.}
The obstruction concerns the bridge, not the elementary localization
mechanism. Let $R>5$, $a_i=R^{2i}$, and define
\[
 M_{ij}=\frac{R^{i+j}}{R^{2i}+R^{2j}}.
\]
Then $M_{ii}=1/2$, and $M_{ij}\leq R^{-|i-j|}$ off the diagonal.
Consequently
\[
 \sum_{j\ne i}M_{ij}<\frac{2}{R-1}<\frac12.                 \tag{7}
\]
For every finite sign word $\sigma$, the vector $M\sigma$ has signs
$\sigma$. Equivalently, the residues $R^j\sigma_j$ in the exterior
Cauchy kernel with queries $a_i$ decode one independent bit per pole.
Reusing these well-separated tokens across arbitrary head slots would
require an additional construction beyond fixed projective changes.

\paragraph{Verification.}
The accompanying Python checker uses only integer and rational
arithmetic. It checks the budget lemma on an exhaustive small grid,
cross-ratios and projective invariance on explicit scales, the robust
error estimate, and strict diagonal dominance for several finite
geometric kernels. These tests supplement the written proof; they do
not establish any growing-head compiler.
\end{document}
